\documentclass[11pt,a4paper]{article}
\usepackage[margin=2.4cm]{geometry}
\usepackage{amsmath,amssymb,bm}
\usepackage{booktabs}
\usepackage{array}
\usepackage{framed}
\usepackage{xcolor}
\usepackage[colorlinks=true,linkcolor=blue,citecolor=blue]{hyperref}

\newcommand{\bO}{\bm{\Omega}}
\newcommand{\bw}{\bm{\omega}}
\newcommand{\bx}{\mathbf{x}}
\newcommand{\bu}{\mathbf{u}}
\newcommand{\bj}{\mathbf{j}}
\newcommand{\bi}{\mathbf{i}}
\newcommand{\bk}{\mathbf{k}}
\newcommand{\bq}{\mathbf{q}}
\newcommand{\bI}{\mathbf{I}}
\newcommand{\It}{\tilde{\mathbf{I}}}
\newcommand{\bG}{\bm{\Gamma}}
\newcommand{\bn}{\hat{\mathbf{n}}}
\newcommand{\E}{\mathsf{E}}
\newcommand{\Id}{\mathsf{1}}
\newcommand{\R}{\mathsf{R}}
\newcommand{\Sm}{\mathsf{S}}
\newcommand{\K}{\mathsf{K}}
\newcommand{\dSL}{\Delta\mathrm{SL}}
\newcommand{\oc}{\mathcal{C}}
\newcommand{\dd}{\,\mathrm{d}}
\newcommand{\pM}{\partial M}

\definecolor{clarbar}{RGB}{70,110,170}
\renewenvironment{leftbar}{%
  \def\FrameCommand{{\color{clarbar}\vrule width 2.5pt}\hspace{8pt}}%
  \MakeFramed{\advance\hsize-\width\FrameRestore}}%
  {\endMakeFramed}
\newenvironment{clar}{\begin{leftbar}\small\noindent\textbf{Clarification.}\ }{\end{leftbar}}

\title{Dynamic rotational feedbacks: solving the forward problem\\[4pt]
\large Companion note --- two solution strategies, and a demonstration}
\author{Working notes (prepared with Claude) --- to be checked}
\date{23 September 2026}

\begin{document}
\maketitle

\tableofcontents

\section{Purpose and scope}

This is a companion to \texttt{rotation\_alattar\_notation.pdf} (hereafter Note~1), which derived the
dynamic rotational theory and its reciprocity properties. Here the concern is narrower and purely
practical: \emph{given a prescribed load history, how do you actually solve the sea-level equation
with a dynamically rotating Earth?}

The immediate aim is a demonstration --- a toy problem in which the Chandler wobble, the
re-orientation of the equatorial bulge, and the resulting redistribution of the oceans can be
computed and animated. That changes what counts as a good method. For an error-budget calculation
one wants the cheapest scheme that gets a small correction right; for a demonstration one wants a
scheme that is cheap to re-run with different parameters, and one wants to choose the forcing so
that the dynamic and static theories visibly disagree.

The note covers:
\begin{itemize}
  \item \S\ref{sec:methodA}: the scheme from Note~1, summarised for forward modelling only (Method~A).
  \item \S\ref{sec:methodB}: an alternative scheme that iterates over the entire time history
  (Method~B), with the conditions it must satisfy to be correct.
  \item \S\ref{sec:compare}: what actually differs between them, and which to use here.
  \item \S\ref{sec:demo}: three demonstration scenarios and what to plot.
\end{itemize}

\paragraph{Notation.} Method~B uses the notation of the Earth-rotation literature (Munk \&
MacDonald 1960), Note~1 that of Al-Attar et al.\ (2024). The correspondence is:

\begin{center}
\small
\begin{tabular}{@{}lll@{}}
\toprule
Quantity & Note~1 / Method~A & Method~B / Munk--MacDonald\\
\midrule
Rotation vector & $\bw_\perp$ (rad\,s$^{-1}$) & $\mathbf{m}=\bw_\perp/\Omega$, complex $m_1+im_2$\\
Excitation & $\bq$, eq.~\eqref{dp1} & $\bm\psi_{\rm geo}=\bq/\Omega$ (usually $\chi$)\\
Inertia perturbation & $\bj=\bi\bO$ & $c_{13},c_{23}$, with $\bj_\perp=\Omega(c_{13},c_{23})$\\
Centrifugal potential & $\psi$ & $\psi_{\rm rot}$\\
Euler frequency & $\sigma_e=\Omega(I_3-I_1)/I_1$, 304.5\,d & $\sigma_e$\\
Chandler frequency & $\sigma_e\sqrt{\det(\Id-\K)}$, an output & $\tilde\sigma_c=\sigma_c+i\alpha$, an input\\
\bottomrule
\end{tabular}
\end{center}

The single most important line of that table is the last one: in Method~A the free wobble frequency
is an \emph{output} of the model, and in Method~B it is an \emph{input}. Much of
\S\ref{sec:compare} follows from that.

\section{Method A: march in time}
\label{sec:methodA}

This section summarises Note~1 \S\S3.5--3.8. Derivations are there; only what is needed to write a
solver is repeated here.

\subsection{The one equation that changes}

Everything in the static theory is unchanged except the rule that fixes the rotation vector. In a
spherically symmetric model, with $I_1=I_2$ and the axial component neglected as usual, that rule
becomes
\begin{equation}
  \boxed{\;\dot\bw_\perp=\sigma_e\,\E\,(\bw_\perp-\bq),\qquad
  \sigma_e=\frac{\Omega(I_3-I_1)}{I_1},\qquad
  \E=\begin{pmatrix}0&-1\\1&0\end{pmatrix}\;}
  \tag{DP1}\label{dp1}
\end{equation}
in place of the static $\bw_\perp=\bq$. Here $\bq=\beta\bm\phi_{21}-\bk_\perp/(I_3-I_1)$ is exactly
the quantity the static solver assigns to $\bw_\perp$: $\beta$ is \texttt{inertia\_factor} and
$\bm\phi_{21}$ the degree-2, order-1 coefficients of the gravitational potential perturbation
excluding $\psi$. The Euler period is $2\pi/\sigma_e=304.5$\,d.

\paragraph{How to read it.} $\E$ rotates by $90^\circ$, so the departure $\bw_\perp-\bq$ is turned
through a right angle and fed back into $\dot\bw_\perp$. The rotation vector therefore
\emph{circles around} $\bq$ instead of sitting on it. For constant $\bq$,
\begin{equation}
  \bw_\perp(t)=\bq+\R(\sigma_e t)\,[\bw_\perp(0)-\bq],
  \tag{DP2}\label{dp2}
\end{equation}
with $\R(\alpha)$ the rotation by $\alpha$: a circle of radius $|\bw_\perp(0)-\bq|$ traced at rate
$\sigma_e$, prograde. The static solution is the centre of that circle.

\subsection{The feedback, and the wobble frequency}

$\bq$ is not external: part of it is generated by $\bw_\perp$ itself, through the deformation raised
by $\psi$ (tidal Love numbers) and the equilibrium ocean pole tide. By linearity
\begin{equation}
  \bq(t)=\bq^{\zeta}(t)+\K\,\bw_\perp(t),\qquad\text{so}\qquad
  \dot\bw_\perp=\sigma_e\E\big[(\Id-\K)\bw_\perp-\bq^\zeta\big],
  \tag{DP3--DP4}\label{dp4}
\end{equation}
where $\bq^\zeta$ is driven by the direct load and $\K$ is a constant $2\times2$ matrix. Without
oceans $\K=\kappa\Id$ with
\begin{equation}
  \kappa=\frac{k_2^T}{k_s}=0.314,\qquad k_s=\frac{3G(I_3-I_1)}{\Omega^2b^5}=0.951 .
  \tag{DP5}\label{dp5}
\end{equation}
The free motion then circles at
\begin{equation}
  \sigma_{\rm CW}=\sigma_e\sqrt{\det(\Id-\K)}\;\xrightarrow[\text{no ocean}]{}\;\sigma_e(1-\kappa),
  \qquad T_{\rm CW}=\frac{304.5\,\mathrm{d}}{1-0.314}\approx444\,\mathrm{d},
  \tag{DP6}\label{dp6}
\end{equation}
against an observed 433\,d. Adding the oceans lengthens this further, by of order tens of days
(Dahlen 1976).

\subsection{Time stepping}

Applying the trapezoidal rule to eq.~\eqref{dp1} with $a=\sigma_e\Delta t/2$,
\begin{equation}
  \boxed{\;\bw_{n+1}=\R\,\bw_n+\Sm\,(\bq_n+\bq_{n+1}),\qquad
  \R=\frac{(1-a^2)\Id+2a\E}{1+a^2},\qquad \Sm=\frac{a^2\Id-a\E}{1+a^2}\;}
  \tag{DP7}\label{dp7}
\end{equation}
$\R$ is exactly a rotation by $2\arctan a$, so the free wobble is carried forward with its amplitude
unchanged and the scheme adds no numerical damping; $\Sm$ is the weight with which the forcing pulls
$\bw_\perp$ towards the static solution over the step. $\R+2\Sm=\Id$, so a steady state with
$\bw=\bq$ is preserved exactly. For monthly steps $a\approx0.31$ and the discrete free period is too
long by 1.5\% (450.6\,d against 444.0\,d). Backward Euler, with $b=\sigma_e\Delta t$, damps the
wobble but collapses onto the static theory as $\Delta t\to\infty$:
\begin{equation}
  \bw_{n+1}=\frac{(\Id+b\E)\,\bw_n+(b^2\Id-b\E)\,\bq_{n+1}}{1+b^2}\;\xrightarrow[b\to\infty]{}\;\bq_{n+1}.
  \tag{DP8}\label{dp8}
\end{equation}

\subsection{Scheme A: modify the existing iteration}

The static solver already iterates, at each time, between the load, the fields and $\bw$. Only the
$\bw$ update changes. Given $\bw_n$, $\bq_n$ and the load $\zeta_{n+1}$:

\begin{center}
\fbox{\begin{minipage}{0.93\linewidth}
\begin{enumerate}
  \item Initialise $\bw^{(0)}=\bw_n$ and $\sigma^{(0)}=\zeta_{n+1}+\rho_w\oc\,\dSL_n$.
  \item For $k=0,1,\dots$ until converged:
  \begin{enumerate}
    \item $\bm\psi_{21}=r\,\bw^{(k)}$, with $r=\sqrt{4\pi/15}\,\Omega b^2$;
    \item $u_{lm}=h_l\sigma^{(k)}_{lm}$ and $\phi_{lm}=k_l\sigma^{(k)}_{lm}$, adding
    $h^T_2\bm\psi_{21}$ and $k^T_2\bm\psi_{21}$ at $(l,m)=(2,1)$;
    \item $\bq^{(k)}=\beta\bm\phi_{21}-\bk_{\perp,n+1}/(I_3-I_1)$;
    \item \textbf{the only change:} $\bw^{(k+1)}=\R\bw_n+\Sm(\bq_n+\bq^{(k)})$, in place of
    $\bw^{(k+1)}=\bq^{(k)}$;
    \item add $\psi$ from $\bw^{(k+1)}$ to $\phi$, form $\dSL$, and update
    $\sigma^{(k+1)}=\zeta_{n+1}+\rho_w\oc\,\dSL$.
  \end{enumerate}
  \item Store $\bw_{n+1}$ and $\bq_{n+1}$ for the next step.
\end{enumerate}
\end{minipage}}
\end{center}

The rotational part of this iteration contracts by at most $\|\Sm\|\,\|\K\|=\kappa a/\sqrt{1+a^2}<\kappa$,
so it converges at least as fast as the static loop, at every forcing frequency.

\subsection{Scheme B: solve the rotation separately}
\label{sec:schemeB}

Because the elastic Earth and the equilibrium ocean respond instantaneously, $\bw$ can be eliminated
exactly:
\begin{enumerate}
  \item \emph{Once:} for $\alpha=1,2$, solve the sea-level equation with no direct load and a
  \emph{prescribed} $\bw_\perp=\mathbf{e}_\alpha$ (impose $\psi$, let the ocean respond, do not
  update $\bw$). This gives the response fields $X^\alpha$ and the columns of $\K$.
  \item \emph{For each time:} solve the sea-level equation for $\zeta(t)$ with
  $\bw_\perp\equiv\mathbf{0}$, i.e.\ rotational feedback switched off. This gives $X^\zeta(t)$ and
  $\bq^\zeta(t)$.
  \item \emph{Then:} integrate eq.~\eqref{dp4} on its own, exactly or with eq.~\eqref{dp7}, whose
  implicit step is the $2\times2$ solve
  $(\Id-\Sm\K)\bw_{n+1}=\R\bw_n+\Sm(\bq_n+\bq^\zeta_{n+1})$ with $\bq_n=\bq^\zeta_n+\K\bw_n$ known.
  \item \emph{Finally:} reassemble $X(t)=X^\zeta(t)+\sum_\alpha\omega_\alpha(t)X^\alpha$.
\end{enumerate}
The total cost is $N_t$ sea-level solves \emph{with rotation switched off}, which are independent of
one another, plus two extra solves. Everything rotational is then a two-dimensional linear ODE.

\subsection{Initial conditions}

From rest, $\bw_0=\mathbf{0}$. From static equilibrium under a pre-existing load, set $\bw_0$ to the
static solution. A free wobble present at the start is not determined by the loads and must be
prescribed.

\subsection{Refinement: retaining the neglected terms}
\label{sec:jdot}

Note~1 drops the $\dot\bj$ terms of the Liouville equation by an ordering argument. They can instead
be kept exactly, at no real cost, because for an elastic Earth with an equilibrium ocean the inertia
perturbation is an \emph{instantaneous} linear functional of the current load and the current
rotation:
\begin{equation}
  \bj_\perp=A\,\mathbf{s}+B\,\bw_\perp,\qquad
  A=(I_3-I_1)\beta k_2,\qquad B=(I_3-I_1)\kappa ,
  \tag{DP9}\label{dp9}
\end{equation}
with $\mathbf{s}$ the degree-2, order-1 coefficients of the load. Then
$\dot\bj_\perp=A\dot{\mathbf{s}}+B\dot\bw_\perp$, and the full equatorial equation becomes
\begin{equation}
  (I_3-I_1-B)\,\bw_\perp+\frac{I_1+B}{\Omega}\,\E\dot\bw_\perp=A\Big(\mathbf{s}-\frac{1}{\Omega}\E\dot{\mathbf{s}}\Big).
  \tag{DP10}\label{dp10}
\end{equation}
This is the same oscillator as eq.~\eqref{dp1} with three changes: the restoring term is softened to
$I_3-I_1-B$, as before; the effective inertia is stiffened to $I_1+B$, because the deformation adds
to the wobble's inertia; and the forcing acquires $-\E\dot{\mathbf{s}}/\Omega$, which is the
classical $\dot c/\Omega$ term of the excitation function (Munk \& MacDonald 1960). The free
frequency is $\sigma_{\rm CW}=\Omega(I_3-I_1-B)/(I_1+B)$.

\paragraph{Integrate the angular momentum, not the rotation.} Eq.~\eqref{dp10} needs
$\dot{\mathbf{s}}$, which is awkward. Stepping the equatorial angular momentum perturbation instead
removes it entirely. With $\mathbf{l}=I_1\bw_\perp+\bj_\perp=(I_1+B)\bw_\perp+A\mathbf{s}$, the
torque-free statement $\dot{\mathbf{l}}=\Omega\E[(I_3-I_1)\bw_\perp-\bj_\perp]$ becomes
\begin{equation}
  \boxed{\;\dot{\mathbf{l}}=\E\Big[\sigma_{\rm CW}\,\mathbf{l}-\frac{\Omega I_3A}{I_1+B}\,\mathbf{s}\Big],
  \qquad \bw_\perp=\frac{\mathbf{l}-A\mathbf{s}}{I_1+B}\;}
  \tag{DP11}\label{dp11}
\end{equation}
a constant-coefficient linear ODE driven by $\mathbf{s}$ alone, with $\bw_\perp$ recovered
algebraically. The integrated quantity is the conserved one, so the scheme cannot drift in angular
momentum.

\paragraph{What it is worth.} $B/I_1=\delta\kappa\approx1.0\times10^{-3}$, so the effective-inertia
correction lengthens the Chandler period by about half a day --- negligible beside the core and
anelasticity uncertainties. The $\dot{\mathbf{s}}$ term is of relative size $\nu/\Omega$, which is
neatly (1 day)/(forcing period): about 3\% at monthly periods and 0.3\% at annual. The one place it
matters qualitatively is a step load: angular momentum is conserved across an instantaneous mass
redistribution, so $\bw_\perp$ must jump immediately to $-A\mathbf{s}/(I_1+B)$ --- the rotation axis
holding its place in space while the figure axis moves under it --- and only then begin to circle.
That jump is about 0.2\% of the eventual static response and of opposite sign, and eq.~\eqref{dp1}
misses it entirely.

\begin{clar}
$A$ and $B$ are pure Love-number expressions only for the solid Earth. If $\mathbf{s}$ is the
\emph{total} load, $B$ is the solid-Earth $\kappa$ but $\mathbf{s}$ still depends on $\bw_\perp$
through the pole tide, so the sea-level iteration remains. If $\mathbf{s}$ is the \emph{direct} load,
$B$ must be the full $\K$ and $A$ must account for the ocean response raised by the direct load;
neither is then a Love-number expression, and both come from the solves of \S\ref{sec:schemeB}. In
that case $B$ is a non-isotropic matrix that does not commute with $\E$, so $\sigma_{\rm CW}$ returns
to the $\sqrt{\det}$ form and the stiffening is anisotropic.

The construction relies on the inertia having no memory. For a viscoelastic Earth $\bj$ is a
convolution over the load history, $\dot\bj$ brings that history with it, and internal variables (or
a Laplace-domain treatment) are needed instead.
\end{clar}

\paragraph{Does it break reciprocity?} No. The criterion, derived in Note~1 \S3.3, is that every
matrix multiplying $\dd/\dd t$ must be a multiple of $\E$, since $\E\,\dd/\dd t$ is symmetric and its
contribution to the forward--adjoint pairing therefore integrates to a boundary term. The $B$ term
satisfies this, simply renormalising $\bG$; the ocean term does so only if the ocean's rotational
response matrix is isotropic, and the violation otherwise is of relative size $\approx2\times10^{-4}$
--- far below the terms being retained here. For an inversion the practical point is unaffected
either way: the adjoint to implement is the transpose of the discrete forward operator, which is
exact regardless of how the continuum statement comes out.

\section{Method B: iterate over the whole time history}
\label{sec:methodB}

Method~A exploits causality: the discretised problem can be marched, so the history is built one step
at a time. The alternative is to treat the entire space--time history as a single unknown and iterate
on it, exactly as the standard pseudo-spectral solver iterates on a single spatial map. This is the
natural approach if the sea-level problem is nonlinear, and it is worth setting out properly rather
than only as a foil for Method~A.

\begin{clar}
The scheme below follows one drafted by Gemini in response to a question from DH. It is restated
here as a method in its own right, with the consistency requirements of \S\ref{sec:consistency}
built into the steps rather than appended as corrections.
\end{clar}

Because the Liouville equation is a differential equation in time, the iteration runs over entire
time histories rather than over spatial grids alone.

\subsection{Assumptions}

\begin{enumerate}
  \item \textbf{Elastic Earth.} The deformation $u$ and gravitational perturbation $\phi$ respond
  instantly to the load; long-term viscoelastic mantle flow is ignored.
  \item \textbf{Hydrostatic oceans.} The Chandler period (about 14 months) is long enough that the
  ocean surface is taken to conform instantly to the prevailing equipotential.
  \item \textbf{No relative angular momentum.} Ice and meltwater act purely as a mass load; they are
  not moving fast enough to contribute a relative angular momentum term ($h_i=0$, equivalently
  $\bk=\mathbf{0}$ in the generalised problem).
  \item \textbf{Damped kinematics.} The wobble is propagated with a damped, complex Chandler
  frequency $\tilde\sigma_c=\sigma_c+i\alpha$ (Munk \& MacDonald 1960), representing anelasticity and
  core coupling.
\end{enumerate}

Assumptions 1--3 are exactly those of Note~1. Assumption~4 adds physics that Note~1 leaves out; it is
a modelling choice rather than a numerical requirement (\S\ref{sec:compare_damping}).

\subsection{The iterative loop}

The unknowns are spatio-temporal fields such as $\dSL(\theta,\lambda,t)$, evaluated on a discretised
time vector $t=[t_0,\dots,t_N]$.

\paragraph{Step 1: initialise.} Guess the sea-level change over the whole history. The simplest
choice is a rigid-Earth ``eustatic'' model in which meltwater spreads uniformly over the oceans at
each time: $\dSL^{(0)}(\theta,\lambda,t)$.

\paragraph{Step 2: build the total load.} At iteration $k$, at every time,
\begin{equation}
  \sigma^{(k)}(t)=\rho_i(1-\oc)\,\Delta I(t)+\rho_w\,\oc\,\dSL^{(k)}(t),
  \tag{B1}\label{b1}
\end{equation}
with $\Delta I(t)$ the prescribed ice history, $\oc$ the ocean function, and the factor $1-\oc$
accounting for floating ice (Note~1, eqs.~2.11--2.12).

\paragraph{Step 3: form the excitation.} At each time, take the degree-2, order-1 coefficients of the
\emph{total} perturbed gravitational potential --- which contains the solid Earth's deformation as
well as the direct mass --- and convert them to the inertia perturbation through MacCullagh's
formula, Note~1 eq.~(P2). In the notation of \S\ref{sec:methodA} this is $\bq$, and the excitation
functions are
\begin{equation}
  \bm\psi^{(k)}_{\rm geo}(t)=\psi^{(k)}_1(t)+i\,\psi^{(k)}_2(t)=\frac{\bq^{(k)}(t)}{\Omega},
  \qquad
  \psi^{(k)}_3(t)=-\frac{c^{(k)}_{33}(t)}{C}.
  \tag{B2}\label{b2}
\end{equation}

\begin{clar}
The deformation contribution must not be omitted. Integrating the surface load alone gives only the
direct mass term; the solid Earth's own redistribution supplies the rest, conventionally as a factor
$(1+k_2')$ with $k_2'\approx-0.30$, i.e.\ a 30\% reduction. Reading the inertia off the total
potential coefficient, as above, includes it automatically and is the safer construction.

The axial relation $\psi_3=-c_{33}/C$ is algebraic rather than an ODE, and correctly so: the axial
Liouville equation integrates exactly, so length of day carries no memory at this order (Note~1,
eq.~D3.5).
\end{clar}

\paragraph{Step 4: propagate the wobble.} With $\bm\psi^{(k)}_{\rm geo}(t)$ known, integrate
\begin{equation}
  \frac{i}{\tilde\sigma_c}\,\dot{\mathbf{m}}^{(k)}(t)+\mathbf{m}^{(k)}(t)=\bm\psi^{(k)}_{\rm geo}(t),
  \tag{B3}\label{b3}
\end{equation}
for the rotation history $\mathbf{m}^{(k)}(t)=\bw^{(k)}_\perp(t)/\Omega$, with an initial condition at
$t_0$ --- typically $\mathbf{m}(0)=0$ for a relaxed state before the melt event. Which feedbacks
$\tilde\sigma_c$ may contain is fixed by \S\ref{sec:consistency}.

\begin{clar}
Eq.~\eqref{b3} is eq.~\eqref{dp1} written complexly, with $\E$ replaced by multiplication by $i$:
rearranging gives $\dot{\mathbf{m}}=i\tilde\sigma_c(\mathbf{m}-\bm\psi_{\rm geo})$. The one
substantive difference is the coefficient: using a Chandler rather than a Euler frequency means part
of the rotational feedback has been folded into the propagator instead of being generated by the
iteration.
\end{clar}

\paragraph{Step 5: centrifugal potential.} From $\mathbf{m}^{(k)}(t)$, and $m_3(t)$ if the axial term
is retained, form $\psi^{(k)}_{\rm rot}(\theta,\lambda,t)$ at every time.

\paragraph{Step 6: deformation and gravity.} Map the load to the response with loading Love numbers
$h_n,k_n$, \emph{and} the centrifugal perturbation to its response with the tidal Love numbers
$h^T_2,k^T_2$:
\begin{equation}
  u_{lm}=h_l\sigma^{(k)}_{lm}+h^T_2\,\psi^{(k)}_{{\rm rot},lm}\big|_{l=2},\qquad
  \phi_{lm}=k_l\sigma^{(k)}_{lm}+k^T_2\,\psi^{(k)}_{{\rm rot},lm}\big|_{l=2}.
  \tag{B5}\label{b5}
\end{equation}

\begin{clar}
Both terms are needed. Folding the feedback into $\tilde\sigma_c$ accounts for the centrifugal
perturbation's effect on the wobble \emph{dynamics}, but the deformation \emph{fields} $u$ and
$\phi$ still require it explicitly. Without the tidal terms the pole-tide deformation of the solid
Earth --- a good part of what a demonstration is meant to show --- is absent from the computed sea
level.
\end{clar}

\paragraph{Step 7: update sea level.} Force the oceans onto the new combined equipotential:
\begin{equation}
  \dSL^{(k+1)}(\theta,\lambda,t)=\frac1g\left(\phi^{(k)}-g\,u^{(k)}+\psi^{(k)}_{\rm rot}\right)+\frac{\Phi_g(t)}{g},
  \tag{B4}\label{b4}
\end{equation}
with $\Phi_g(t)$ fixed at each time by mass conservation between ocean and ice.

\begin{clar}
Eq.~\eqref{b4} uses the opposite sign convention for $\phi$ and $\psi$ from Al-Attar et al., who
write $\dSL=-(gu+\phi+\psi)/g+\Phi_g/g$ with $\phi$ the Newtonian potential perturbation, so that the
geoid is $-\phi/g$. Eq.~\eqref{b4} takes the geoid as $+\phi/g$ and flips $\psi$ to match. Each is
self-consistent; mixing them is fatal, so settle on one before implementing anything from both
documents.
\end{clar}

\paragraph{Step 8: iterate.} Compare $\dSL^{(k+1)}$ with $\dSL^{(k)}$ over all grid points and all
times; if the difference exceeds tolerance, return to Step~2.

Convergence is a global condition because the coupling runs both ways in time: the wobble at
$t_{100}$ depends on the sea level at $t_{50}$ through the excitation history, and updating the sea
level in Step~7 changes the whole wobble history, hence the whole centrifugal history.

\subsection{Consistency requirements}
\label{sec:consistency}

Three conditions must hold for the loop to converge to the right answer. Each is a case where the
method returns a plausible but wrong result rather than failing outright.

\begin{enumerate}
  \item \textbf{The inertia must include the deformation} (Step~3). Otherwise the excitation is too
  large by about 30\%.
  \item \textbf{The fields must include the tidal response} (Step~6), or the solid Earth's pole-tide
  deformation is missing from $u$, $\phi$ and hence from sea level.
  \item \textbf{Each feedback must be counted exactly once} (Step~4). The rotational feedback has two
  parts: the solid Earth's elastic yielding to the centrifugal perturbation (the factor $1-\kappa$ of
  eq.~\eqref{dp6}), and the equilibrium ocean pole tide. Either may be folded into $\tilde\sigma_c$ or
  generated by the outer loop, but not both. The two self-consistent choices are:
  \begin{itemize}
    \item $\tilde\sigma_c$ = the \emph{no-ocean} Chandler frequency, with the outer loop generating
    the ocean pole tide through Step~7. This is the recommended choice.
    \item $\tilde\sigma_c=\sigma_e$, the Euler frequency, with the outer loop generating both parts.
    Simplest to state, but with much the worse convergence (\S\ref{sec:gain}).
  \end{itemize}
  Using the \emph{observed} 433\,d value is neither: it already contains the ocean pole tide, which
  Step~7 then adds a second time.
\end{enumerate}

\section{What actually differs, and which to use}
\label{sec:compare}

\subsection{The two methods solve the same equations}

This is worth stating plainly before the differences: eq.~\eqref{b3} is eq.~\eqref{dp1}, and
eq.~\eqref{b4} is the sea-level equation with the centrifugal term. Where both converge, they
converge to the same answer. The differences below are about cost, robustness, and which quantities
are inputs rather than outputs, provided the consistency requirements of \S\ref{sec:consistency}
are met.

\subsection{Global iteration versus causal marching}

Method~B iterates over the whole time history: guess $\dSL(\theta,\lambda,t)$ everywhere, rebuild
the excitation over all $t$, re-integrate the ODE over all $t$, repeat. Method~A marches, because
the discretised problem is causal: at step $n+1$ the only things needed from the past are the
2-vectors $\bw_n$ and $\bq_n$. There is no outer loop over time, and no need to hold the entire
space--time field in memory.

For a forward run this is the difference between $N_{\rm iter}\times N_t$ sea-level solves, held
together by a global convergence test, and $N_t$ solves each converging locally --- or, with
Scheme~B, $N_t$ \emph{independent} solves that can be run in any order or in parallel.

\subsection{Iterating the feedback versus inverting it}
\label{sec:gain}

Method~B iterates the rotational feedback; Method~A inverts it. In Method~B's outer loop, one lap of
the cycle (sea level $\to$ load $\to$ excitation $\to$ $\mathbf{m}$ $\to$ $\psi_{\rm rot}$ $\to$ sea
level) has gain
\begin{equation}
  A(\nu)=\frac{\kappa_{\rm fb}}{\left|1-\nu/\sigma_c\right|}
  \tag{C1}\label{c1}
\end{equation}
at forcing frequency $\nu$, where $\kappa_{\rm fb}$ is whatever part of the feedback is not already
inside $\sigma_c$. The iteration diverges where $A>1$. Two cases:
\begin{itemize}
  \item With $\tilde\sigma_c$ the no-ocean Chandler frequency, the solid-Earth part is already folded
  in and $\kappa_{\rm fb}$ is only the ocean pole tide, of order $0.05$. The loop diverges for forcing
  within roughly 5\% of the Chandler period, say 425--465\,d, and converges slowly either side.
  \item With $\tilde\sigma_c=\sigma_e$, the Euler frequency, the loop must generate both parts, so
  $\kappa_{\rm fb}=\kappa=0.314$ and the divergent band widens to roughly 230--444\,d --- which
  swallows the annual band. This is the price of the simpler propagator.
\end{itemize}
Method~A has no such band. In Scheme~B the $2\times2$ system $(\Id-\Sm\K)\bw_{n+1}=\dots$ is solved
directly, and in Scheme~A the per-step contraction is $\kappa a/\sqrt{1+a^2}<\kappa$ at every
frequency, because the step map $\Sm$ carries a factor $a/\sqrt{1+a^2}<1$.

For the demonstration this matters more than it would for a production run, because the scenarios
worth showing (\S\ref{sec:demo}) are precisely the ones near resonance.

\subsection{Damping}
\label{sec:compare_damping}

Method~B propagates the wobble with a complex $\tilde\sigma_c$, which is partly motivated by not
wanting the wobble to ring indefinitely in the solver. Two comments. The ringing is physical, not
numerical: an undamped elastic Earth
genuinely rings, and eq.~\eqref{dp7} reproduces that faithfully because $\R$ is an exact rotation.
And damping is worth including --- the observed $Q_{\rm CW}$ is of order $10^2$ (Smith \& Dahlen 1981) --- but as a physical
statement about anelasticity and core coupling, not as a device to stabilise an integrator. For the
demonstration it is better treated as an adjustable knob (\S\ref{sec:demo_impl}).

\subsection{What each method is good for}

Method~A, for this problem:
\begin{itemize}
  \item \emph{Cheap animation.} With Scheme~B, once $\K$ and $X^\alpha$ exist, rendering the wobble
  at any frame rate costs a $2\times2$ ODE and no sea-level solves.
  \item \emph{Cheap parameter sweeps.} Changing the damping, retuning $\sigma_e$ for core
  decoupling, or switching rotation off entirely costs nothing further. Every such experiment in
  Method~B is a complete re-run.
  \item \emph{$\sigma_{\rm CW}$ is diagnosed, not imposed,} so it can be checked against 433\,d, and
  double counting cannot hide.
  \item \emph{No artificial damping needed} for stability.
\end{itemize}

Method~B, in its favour:
\begin{itemize}
  \item \emph{Nonlinearity.} A global Picard loop does not care whether the problem is linear. With
  moving shorelines, grounded-ice margins or a viscoelastic Earth, Scheme~B's exact superposition
  fails (Scheme~A largely survives, being a fixed point per step). Method~B's structure generalises
  unchanged.
  \item \emph{Contact with the rotation literature.} Excitation functions $\chi$ are how observed and
  modelled excitation budgets are published, so working in those variables saves a conversion when
  comparing against them. This is a change of variables ($\mathbf{m}=\bw_\perp/\Omega$,
  $\bm\chi=\bq/\Omega$), not a property of the scheme, and is equally available to Method~A --- a
  convenience rather than a reason to choose one method over the other.
\end{itemize}

\subsection{Recommendation for the demonstration}

Use Method~A in the Scheme~B form. The problem is linear and time-invariant, which is exactly the
case where the rotational part collapses to two extra solves plus a two-dimensional ODE, and the
demonstration needs many re-runs with different forcing periods and damping. Keep Method~B's
structure in reserve if the toy problem later grows a moving shoreline. In either case, meet the
consistency requirements of \S\ref{sec:consistency} and settle the sign convention of
eq.~\eqref{b4} first.

\section{Designing the demonstration}
\label{sec:demo}

For an error budget one asks whether the dynamic correction is above the noise; the answer, for
multi-year trends, is no. For a demonstration the criterion is the opposite: pick the regimes where
the dynamic theory does something the static theory \emph{cannot do at all}. There are three.

\subsection{Scenario 1: impulsive melt, then nothing}

Apply a load over a month, then hold everything fixed. The static theory gives a sea-level field
that then sits still forever. The dynamic theory gives a rotation vector that circles its new
equilibrium at the Chandler rate indefinitely, so the sea-level field keeps moving for years after
the forcing has stopped.

This is the strongest demonstration, because the difference is qualitative rather than a
percentage: ringing after the load stops is pure memory, and there is nowhere for it to come from in
the static formulation. A step load is also the one case where the refinement of \S\ref{sec:jdot} is
worth carrying: it supplies the instantaneous jump in $\bw_\perp$ demanded by conservation of angular
momentum at the moment the mass is added.

\subsection{Scenario 2: periodic forcing swept through resonance}

Melt and refreeze periodically, sweeping the period from about 100\,d to about 2000\,d. Plot the
amplitude of the rotational part of the sea-level signal against forcing period. The static theory
is a flat line; the dynamic theory is a resonance curve peaking at $T_{\rm CW}$, with a sign
reversal either side. From eq.~\eqref{dp4} with $\K=\kappa\Id$, the ratio of the two is
\begin{equation}
  \frac{w_{\rm dynamic}}{w_{\rm static}}=\frac{1}{1-\nu/\sigma_{\rm CW}},
  \tag{C2}\label{c2}
\end{equation}
which is $-0.072$ at 30\,d, $-4.64$ at one year (prograde) and $1$ for a trend. This single plot
conveys the whole content of the extension.

\subsection{Scenario 3: steady trend}

The reassuring limit. The two theories agree on the trend; the dynamic solution differs only by a
constant offset perpendicular to it, of magnitude (trend rate) $\times\,T_{\rm CW}/2\pi\approx70$\,d,
plus a free wobble excited at switch-on.

\subsection{What to plot}

The rotational part of the field is exactly degree 2, order 1, so it is a clean four-lobed pattern
with extrema near $\pm45^\circ$ latitude, and the wobble carries it round in longitude with the
Chandler period. Three panels side by side work well:
\begin{itemize}
  \item \textbf{the pole path}: $\omega_1$ against $\omega_2$, showing the spiral circling the moving
  static solution $\bq$ --- this is the polar motion itself, and eq.~\eqref{dp2} made visible;
  \item \textbf{sea level}: the difference between the rotating and non-rotating solutions, i.e.\
  $\sum_\alpha\omega_\alpha(t)X^\alpha$ from Scheme~B, with the pattern rotating in longitude;
  \item \textbf{the solid Earth}: vertical displacement, showing the equatorial bulge re-orienting.
\end{itemize}
Worth saying in the narration: in an elastic Earth the bulge follows the rotation axis
\emph{instantaneously}. There is no lag between axis and bulge; the lag is between the axis and the
forcing. A lagging bulge requires anelasticity, which is also what damps the wobble.

\subsection{Implementation notes}
\label{sec:demo_impl}

\begin{itemize}
  \item Use Scheme~B. Two solves give $X^1,X^2$ and $\K$; one solve per time gives the non-rotating
  background; after that the animation is a $2\times2$ ODE, so frame rate and forcing period are
  free parameters.
  \item $l_{\max}=64$ is ample: the rotational signal is degree 2, order 1 exactly, and the
  background fingerprint does not need to be sharp for this purpose.
  \item Make damping a slider rather than a fixed choice. Showing the same melt event with
  $Q=\infty$ and $Q=100$ makes the point that the ringing is physical and its decay is a modelling
  assumption.
  \item Diagnose $\sigma_{\rm CW}$ from $\K$ and display it. Running once with the ocean function
  zeroed and once with real oceans shows the period lengthening by of order tens of days purely
  because of the water --- a result worth displaying in its own right, and the cleanest way to show that the
  feedback is counted once (\S\ref{sec:consistency}, item~3).
\end{itemize}

\subsection{Two cautions for the narration}

\paragraph{``Sloshing'' here means equilibrium redistribution.} In this theory the ocean surface is
always exactly an equipotential and the water has no inertia of its own. What moves is the
equipotential; the water follows it instantly. All the memory in the system lives in the rotation
axis. Genuine sloshing --- waves, adjustment, a non-equilibrium pole tide --- would need a
shallow-water model on top of this. The animation will look like a dynamic ocean response, and it is
not one, so this is worth stating explicitly.

\paragraph{Amplitudes.} A realistic melt event produces polar motion of order milliarcseconds and a
pole tide of order hundredths of a millimetre, which will be invisible next to the direct
fingerprint. For the demonstration, either exaggerate the load heavily or --- better --- plot the
rotational component on its own colour scale with the direct part subtracted, which is both more
honest and shows the pattern more clearly. For scale, the equilibrium pole tide associated with the
observed Chandler wobble is of order a few millimetres.

\begin{thebibliography}{9}
\bibitem{aa24} Al-Attar, D., Syvret, F., Crawford, O., Mitrovica, J.X.\ \& Lloyd, A.J., 2024. Reciprocity and sensitivity kernels for sea level fingerprints, \emph{Geophys.\ J.\ Int.}, 236(1), 362--378.
\bibitem{dahlen76} Dahlen, F.A., 1976. The passive influence of the oceans upon the rotation of the Earth, \emph{Geophys.\ J.\ R.\ astr.\ Soc.}, 46, 363--406.
\bibitem{munk} Munk, W.H.\ \& MacDonald, G.J.F., 1960. \emph{The Rotation of the Earth}, Cambridge University Press.
\bibitem{sd81} Smith, M.L.\ \& Dahlen, F.A., 1981. The period and Q of the Chandler wobble, \emph{Geophys.\ J.\ R.\ astr.\ Soc.}, 64, 223--281.
\end{thebibliography}

\end{document}
